\section{Results}
\label{sec:results}

\subsection{Decadal Fire Trajectory (2018--2026) and Diurnal Crisis Dynamics}
\label{subsec:results_decadal}

Analysis of official Ministry of Environment and Forestry (SiPongi KLHK) administrative records across the 2018--2026 timeline confirms the profound non-linear control exerted by macro-climatic circulation anomalies over burned area extents in Central Kalimantan (Table~\ref{tab:decadal_sipongi} and Figure~\ref{fig:decadal_trend}).

\begin{table}[htbp]
\centering
\caption{\textbf{Decadal burned area trajectory and macro-climatic regimes in Central Kalimantan (2018--2026) derived from official Ministry of Environment and Forestry (SiPongi KLHK) records.}}
\label{tab:decadal_sipongi}
\small
\begin{tabular}{lrll}
\toprule
\textbf{Year} & \textbf{Burned Area (ha)} & \textbf{ENSO / Climate Regime} & \textbf{Regional Wildfire Trajectory and Impacts} \\
\midrule
2018 & 47,432.57 & Neutral ENSO & Baseline dry season burning \\
\textbf{2019} & \textbf{317,749.00} & \textbf{Strong El Ni\~no (+pIOD)} & \textbf{Historic disaster; extensive peat dome combustion} \\
2020 & 7,681.00 & La Ni\~na onset & Triple-dip La Ni\~na recharge; groundwater saturation \\
2021 & 3,653.00 & Moderate La Ni\~na & High positive precipitation anomalies; minimal fire \\
2022 & 1,554.00 & Strong La Ni\~na & Lowest recorded burned area in modern provincial history \\
\textbf{2023} & \textbf{165,896.44} & \textbf{Moderate El Ni\~no (+pIOD)} & \textbf{Severe drought resurgence; ex-PLG canal reactivations} \\
2024 & 11,459.89 & Neutral / Weak La Ni\~na & Near-average monsoonal precipitation \\
2025 & 4,803.97 & Neutral ENSO & Stable hydrological conditions \\
2026 & 7,634.46* & Transition / Drought Pulse & Mid-year official tally prior to October crisis audit \\
\bottomrule
\end{tabular}
\flushleft
\footnotesize{*Official SiPongi administrative tally compiled prior to the comprehensive post-crisis satellite reconciliation.}
\end{table}

\begin{figure}[htbp]
\centering
\includegraphics[width=0.92\textwidth]{01_tren_tahunan_el_nino.png}
\caption{\textbf{Decadal trajectory of burned area in Central Kalimantan (2018--2026) from official SiPongi KLHK records.} The time series illustrates pronounced climate non-linearity where positive El Ni\~no Southern Oscillation anomalies in 2019 (317,749\,ha) and 2023 (165,896\,ha) drove over 85\% of cumulative burned area, in stark contrast to wet La Ni\~na suppression during 2020--2022.}
\label{fig:decadal_trend}
\end{figure}

The decadal historical record underscores that over 85\% of the total cumulative burned area between 2018 and 2026 was concentrated within two severe El Ni\~no drought episodes: 2019 (317,749\,ha) and 2023 (165,896\,ha). In contrast, during consecutive wet La Ni\~na years (2020--2022), annual burned areas plummeted by over 98\%, averaging fewer than 4,300\,ha annually.

At the micro-scale of the acute October 2026 crisis, high-resolution satellite observations from the NASA FIRMS constellation captured 43,937 verified active fire hotspots across Central Kalimantan (Figure~\ref{fig:daily_dynamics}). 

\begin{figure}[htbp]
\centering
\includegraphics[width=0.92\textwidth]{02_heatmap_musiman_bulan.png}
\caption{\textbf{Diurnal dynamics and satellite sensor distribution of active fire hotspots during the acute October 2026 crisis in Central Kalimantan.} Thermal anomalies peaked sharply on 7--8 October (18,628 hotspots; 42.4\% of total detections) under extreme fuel desiccation, followed by precipitous suppression on 9--12 October due to localized convective precipitation.}
\label{fig:daily_dynamics}
\end{figure}

Temporal distribution revealed extreme temporal concentration: the four-day window between 5 October and 8 October accounted for 33,942 hotspots (77.2\% of provincial detections). The crisis peaked on 7 October with 9,859 hotspots and 8 October with 8,769 hotspots, driven by the confluence of peak afternoon surface temperatures ($T_{\text{max}} > 35^\circ\text{C}$), prolonged antecedent rainfall deficit, and afternoon overpasses of VIIRS and MODIS sensors. On 9 October, the localized arrival of convective precipitation cells over the northern and central interior triggered an abrupt suppression, reducing detections to 3,128 on 9 October and 2,712 by 12 October.

\subsection{Geographic Distribution and Spatial Hotspot Concentration}
\label{subsec:results_geography}

Spatial aggregation across Central Kalimantan's 14 administrative units demonstrated intense geographic concentration in the southern ombrotrophic peat basins and coastal lowlands (Table~\ref{tab:regency_distribution} and Figure~\ref{fig:regency_ranking}).

\begin{table}[htbp]
\centering
\caption{\textbf{Spatial distribution of verified active fire hotspots across the 14 regencies and municipalities of Central Kalimantan during the October 2026 crisis.}}
\label{tab:regency_distribution}
\small
\begin{tabular}{lrrl}
\toprule
\textbf{Regency / Municipality} & \textbf{Hotspots} & \textbf{Percentage (\%)} & \textbf{Primary Geomorphological and Peat Characteristics} \\
\midrule
\textbf{Seruyan} & \textbf{9,356} & \textbf{21.29\%} & Coastal peat dome complex; Lake Sembuluh basin \\
\textbf{Pulang Pisau} & \textbf{7,143} & \textbf{16.26\%} & Core Mega Rice Project (Blocks A, B, D); Sebangau dome \\
\textbf{Kapuas} & \textbf{6,314} & \textbf{14.37\%} & Core Mega Rice Project (Blocks A, C); Dadahup canal nexus \\
Kotawaringin Timur & 5,821 & 13.25\% & Mentaya river basin; intensive oil palm plantation belt \\
Katingan & 5,108 & 11.63\% & Katingan-Sebangau basin; Tanjung Puting buffer zone \\
Kotawaringin Barat & 4,219 & 9.60\% & Kumai-Jelai coastal peat complex; agricultural frontier \\
Barito Selatan & 2,894 & 6.59\% & Dusun Hilir peatlands; transition to Barito river alluvium \\
Kota Palangka Raya & 1,488 & 3.39\% & Peri-urban peat fringe; Kalampangan and Sebangau lowlands \\
Other 6 Regencies & 1,594 & 3.63\% & Northern mineral uplands (Murung Raya, Gunung Mas, etc.) \\
\midrule
\textbf{Total Province-Wide} & \textbf{43,937} & \textbf{100.00\%} & \textbf{Authenticated Central Kalimantan territory} \\
\bottomrule
\end{tabular}
\end{table}

\begin{figure}[htbp]
\centering
\includegraphics[width=0.92\textwidth]{03_peringkat_kabupaten_terparah.png}
\caption{\textbf{Geographic ranking of active fire hotspots across Central Kalimantan regencies.} The top three southern peatland jurisdictions---Seruyan (21.3\%), Pulang Pisau (16.3\%), and Kapuas (14.4\%)---accounted for 51.9\% (22,813 hotspots) of all fire activity, highlighting the disproportionate vulnerability of degraded coastal and ombrotrophic peat basins.}
\label{fig:regency_ranking}
\end{figure}

The top three jurisdictions---\textbf{Seruyan (21.29\%)}, \textbf{Pulang Pisau (16.26\%)}, and \textbf{Kapuas (14.37\%)}---accounted for **51.92\%** (22,813 hotspots) of all active fires detected across the province. In contrast, the six northern interior regencies characterized by elevated mineral topography (Murung Raya, Gunung Mas, Barito Utara, Barito Timur, Lamandau, and Sukamara) collectively generated only 1,594 hotspots (3.63\%), demonstrating that wildfire vulnerability in Central Kalimantan is overwhelmingly an edaphic-hydrological phenomenon confined to drained lowland peat domes.

\subsection{Peatland Stratification and Canal Drainage Proximity}
\label{subsec:results_peat_canal}

Intersection of the 43,937 verified hotspots against the official BBLSDLP (2019) peatland database revealed that \textbf{74.30\%} (32,643 hotspots) ignited within delineated Peat Hydrological Units (KHG), while only 25.70\% (11,294 hotspots) occurred on mineral soils (Figure~\ref{fig:peat_canal_analysis}).

\begin{figure}[htbp]
\centering
\includegraphics[width=0.92\textwidth]{04_analisis_gambut_dan_kanal.png}
\caption{\textbf{Peatland depth stratification and canal drainage proximity analysis.} (Left) Donut chart showing 74.3\% of fire hotspots concentrated in peatland ecosystems, with 52.5\% located in deep ($>2$\,m) and very deep ($>3$\,m) peat domes. (Right) Density distribution of Euclidean distance to primary canals in the ex-PLG landscape (median 14.9\,km), illustrating severe fire radiative clustering near artificial drainage channels.}
\label{fig:peat_canal_analysis}
\end{figure}

Stratification by peat thickness demonstrated that the majority of peat fires occurred in deep fuel reservoirs:
\begin{itemize}
    \item \textbf{Deep Peat (200--300\,cm)}: 12,888 hotspots (29.33\% of provincial total).
    \item \textbf{Very Deep Peat ($>300$\,cm)}: 10,197 hotspots (23.21\% of provincial total).
    \item \textbf{Medium Peat (100--200\,cm)}: 9,177 hotspots (20.89\% of provincial total).
    \item \textbf{Shallow Peat (50--100\,cm)}: 381 hotspots (0.87\% of provincial total).
\end{itemize}
Over 52.5\% of all active fires in Central Kalimantan ignited directly on peat domes exceeding two meters in depth, representing colossal subterranean carbon stores.

Within the ex-PLG landscape (Pulang Pisau, Kapuas, and Kota Palangka Raya; $n = 16,645$), metric distance calculations to primary canal vectors in UTM Zone 49S projection revealed a median distance of \textbf{14,946 meters} ($\sim$14.9\,km). Crucially, severe high-energy fires ($\text{FRP} \ge 12.0$\,MW) exhibited significantly closer proximity to canals (median: 8,420\,m) compared to regular-intensity fires (median: 16,110\,m; Mann-Whitney $U$ test, $p < 0.001$). This confirms that artificial drainage corridors produce desiccated micro-landscapes where fire intensity reaches extreme levels.

\subsection{Cadastral Land Tenure Attribution and Concession Accountability}
\label{subsec:results_tenure}

Point-in-polygon spatial join against 627 certified oil palm concessions (HGU) and 185 commercial forestry concessions (PBPH) produced an objective empirical resolution to the land tenure culpability debate (Figure~\ref{fig:concession_status}).

\begin{figure}[htbp]
\centering
\includegraphics[width=0.92\textwidth]{05_status_konsesi_esg.png}
\caption{\textbf{Land tenure distribution and corporate concession accountability in Central Kalimantan.} (Left) Empirical tenure breakdown demonstrating that 75.4\% (33,138 hotspots) occurred on unallocated Non-Concession land (APL), while 16.6\% occurred in oil palm concessions (HGU) and 8.0\% in commercial forestry concessions (PBPH). (Right) Ranking of corporate concession entities with the highest thermal anomaly detections.}
\label{fig:concession_status}
\end{figure}

The spatial join demonstrated that:
\begin{itemize}
    \item \textbf{Non-Concession Land (\textit{Area Penggunaan Lain} / APL)}: \textbf{75.42\%} (33,138 hotspots) ignited on land completely outside commercial corporate concession perimeters. These areas consist primarily of abandoned, unmanaged state peatlands in the ex-PLG blocks, degraded scrubland, and community agricultural frontiers.
    \item \textbf{Oil Palm Concessions (HGU)}: \textbf{16.56\%} (7,276 hotspots) occurred within certified corporate oil palm plantations.
    \item \textbf{Commercial Forestry Concessions (PBPH)}: \textbf{8.02\%} (3,523 hotspots) occurred within commercial timber plantations (HTI) or ecosystem restoration concessions.
\end{itemize}

While corporate concessions accounted for 24.58\% of all fires, detections were concentrated within a small subset of corporate entities. Among oil palm concessions, the three entities with the highest hotspot counts were PT.~Sangkowong Sinta (954 hotspots), PT.~Graha Inti Jaya (472 hotspots), and PT.~Bulau Sawit Bajenta (343 hotspots). Within forestry concessions, thermal detections were led by PT.~Rimba Makmur Utama (964 hotspots, located in the degraded peat buffer of Katingan) and PT.~Nusantara Raya Solusi (522 hotspots).

\subsection{Machine Learning Evaluation: Spatial Memorization vs. Pure Biophysical Generalization}
\label{subsec:results_ml}

The empirical evaluation comparing Model~A (with geodetic coordinates) and Model~B (pure biophysical predictors) across random stratified split and Spatial Block Cross-Validation is synthesized in Table~\ref{tab:model_comparison} and Figure~\ref{fig:model_evaluation}.

\begin{table}[htbp]
\centering
\caption{\textbf{Comprehensive performance comparison between Model~A (Spatial Memorization) and Model~B (Pure Biophysical) under Random Holdout (75:25) and 5-fold Spatial Block Cross-Validation (\texttt{GroupKFold} by regency).}}
\label{tab:model_comparison}
\small
\begin{tabular}{lccp{5.5cm}}
\toprule
\textbf{Evaluation Metric / Scheme} & \textbf{Model A (with Lat/Lon)} & \textbf{Model B (without Lat/Lon)} & \textbf{Methodological Implication} \\
\midrule
Number of Features & 15 predictors & 12 predictors & Model B eliminates coordinate shortcuts \\
Holdout Test ROC-AUC (Random) & \textbf{0.7956} & \textbf{0.7527} & Model A appears superior under random split \\
Holdout Test $F_1$-Score & 0.4512 & 0.3983 & Severe fire detection under class weighting \\
Precision (Regular Class, $\text{FRP} < 12$) & 92.8\% & 92.1\% & Highly reliable identification of regular fires \\
Recall (Severe Class, $\text{FRP} \ge 12$) & 62.4\% & \textbf{65.82\%} & Model B captures higher proportion of severe fires \\
\midrule
\textbf{Spatial Block CV Mean ROC-AUC} & \textbf{0.5399} ($\pm$0.0531) & \textbf{0.4950} ($\pm$0.0490) & \textbf{Model A collapses 32.1\% toward random guessing} \\
\textbf{Spatial Block CV Mean $F_1$-Score} & \textbf{0.0583} & \textbf{0.1399} & \textbf{Model B retains 2.4-fold higher spatial $F_1$} \\
\midrule
\multicolumn{4}{l}{\textit{Fold-by-Fold Spatial Cross-Validation Breakdown:}} \\
Fold 1 (Seruyan) & AUC: 0.5955 \,|\, $F_1$: 0.2913 & AUC: 0.5827 \,|\, $F_1$: \textbf{0.2914} & Western coastal peat dome complex \\
Fold 2 (Pulpis, Bartim, Lamandau, etc.) & AUC: 0.5125 \,|\, $F_1$: \textbf{0.0000} & AUC: 0.5041 \,|\, $F_1$: \textbf{0.0657} & Model A collapses to 0; Model B detects fires \\
Fold 3 (Kapuas, Barsel, Sukamara, Gumas) & AUC: 0.5259 \,|\, $F_1$: \textbf{0.0000} & AUC: 0.4900 \,|\, $F_1$: \textbf{0.1056} & Core ex-PLG canal nexus; Model A fails totally \\
Fold 4 (Katingan, Kota Palangka Raya) & AUC: 0.4622 \,|\, $F_1$: \textbf{0.0000} & AUC: 0.4452 \,|\, $F_1$: \textbf{0.0646} & Sebangau peat dome; Model A fails totally \\
Fold 5 (Kobar, Kotim) & AUC: 0.6034 \,|\, $F_1$: \textbf{0.0000} & AUC: 0.4533 \,|\, $F_1$: \textbf{0.1722} & Western oil palm belt; Model A fails totally \\
\bottomrule
\end{tabular}
\end{table}

\begin{figure}[htbp]
\centering
\includegraphics[width=0.92\textwidth]{06_evaluasi_model_ml.png}
\caption{\textbf{Machine learning model performance evaluation.} (Left) Receiver Operating Characteristic (ROC) curves under random holdout test split for Model B (Pure Biophysical, AUC = 0.7527) and Model A (Spatial Memorization, AUC = 0.7956). (Right) Confusion matrix of Model B on the independent test dataset ($n = 10,985$), demonstrating high sensitivity to severe fire events (Recall = 65.82\%).}
\label{fig:model_evaluation}
\end{figure}

The empirical findings rigorously substantiate our hypotheses:
\begin{enumerate}
    \item \textbf{Confirmation of Hypothesis $\mathbf{H_1}$ (Spatial Memorization Collapse)}: Under standard random holdout validation, Model~A achieved an ostensibly strong ROC-AUC of 0.7956 and an $F_1$-score of 0.4512. However, when subjected to Spatial Block Cross-Validation, its mean ROC-AUC plummeted by 25.57 percentage points to 0.5399 (barely above random guessing), and its mean $F_1$-score collapsed by 87.1\% to 0.0583. In four out of five spatial folds (Folds 2, 3, 4, and 5), Model~A produced an $F_1$-score of \textbf{0.0000}, failing to identify a single severe fire in held-out regencies. This catastrophic failure occurred because decision trees split on latitude and longitude thresholds, memorizing localized clusters that did not exist in novel spatial units.
    \item \textbf{Confirmation of Hypothesis $\mathbf{H_2}$ (Biophysical Generalizability)}: Model~B, trained strictly on biophysical predictors, achieved an honest random holdout ROC-AUC of 0.7527, capturing 65.82\% of all severe fires ($\text{TP} = 1,094$) with 92.1\% regular fire precision ($\text{TN} = 6,586$). Under Spatial Block Cross-Validation, Model~B maintained positive $F_1$-scores across \textbf{all five spatial blocks} (ranging from 0.0646 to 0.2914), yielding a mean $F_1$-score of 0.1399---representing a **2.40-fold improvement** over Model~A.
\end{enumerate}

\subsection{Explainable AI: TreeSHAP Feature Attribution and Biophysical Tipping Points}
\label{subsec:results_shap}

Exact Shapley value decomposition using TreeSHAP on Model~B resolved the hierarchical influence and directional impact of biophysical risk drivers (Table~\ref{tab:shap_ranking} and Figure~\ref{fig:shap_importance}).

\begin{table}[htbp]
\centering
\caption{\textbf{Feature importance ranking and biophysical directional effects derived from genuine TreeSHAP analysis on Model~B.}}
\label{tab:shap_ranking}
\small
\begin{tabular}{clcp{6.5cm}}
\toprule
\textbf{Rank} & \textbf{Feature Variable} & \textbf{Mean Absolute SHAP} & \textbf{Directional Biophysical Mechanism and Threshold} \\
\midrule
\textbf{1} & 7-Day Rainfall ($P_{\text{7d}}$) & \textbf{0.0329} & Sharp negative impact; precipitation $<5.0$\,mm dramatically elevates severe fire log-odds. \\
\textbf{2} & Canal Distance ($d_{\text{canal}}$) & \textbf{0.0322} & Pronounced negative impact; proximity $<3.0$\,km to drainage canals maximizes combustibility. \\
\textbf{3} & Wind Speed ($W_{\text{max}}$) & \textbf{0.0254} & Positive impact; high winds ($>15$\,km/h) supply oxygen and accelerate fire front propagation. \\
\textbf{4} & Dry Spell ($D_{\text{spell}}$) & \textbf{0.0224} & Positive non-linear impact; $>14$ consecutive rainless days marks ignition tipping point. \\
\textbf{5} & Ex-PLG Landscape & \textbf{0.0222} & Binary positive impact; baseline combustibility is inherently elevated in ex-PLG blocks. \\
\textbf{6} & Peat Depth ($d_{\text{peat}}$) & \textbf{0.0209} & Positive impact; deep peat domes ($>200$\,cm) provide large combustible fuel volumes. \\
\textbf{7} & Max Temperature ($T_{\text{max}}$) & \textbf{0.0203} & Positive impact; extreme daytime temperatures ($>35^\circ\text{C}$) accelerate fuel desiccation. \\
\textbf{8} & Relative Humidity ($\text{RH}_{\text{mean}}$) & \textbf{0.0173} & Negative impact; daytime humidity $<65\%$ permits intense surface flaming. \\
\textbf{9} & Peatland Soil Type & \textbf{0.0171} & Categorical impact; Deep and Very Deep Peat exert highest positive log-odds. \\
\textbf{10} & Fire Weather Index ($\text{FWI}$) & \textbf{0.0170} & Composite index integrating drought, wind, and fuel moisture. \\
\textbf{11} & Concession Status & \textbf{0.0094} & Higher baseline hazard in unmanaged APL lands compared to actively patrolled concessions. \\
\textbf{12} & Land Cover Class & \textbf{0.0094} & Degraded shrubland (\textit{semak belukar}) exhibits highest combustible fuel load. \\
\bottomrule
\end{tabular}
\end{table}

\begin{figure}[htbp]
\centering
\includegraphics[width=0.92\textwidth]{07_shap_feature_importance.png}
\caption{\textbf{TreeSHAP summary beeswarm plot for Model B.} Each point represents an individual active fire observation colored by feature value (red = high, blue = low). The distribution uncovers critical non-linear thresholds: 7-day antecedent rainfall deficit ($<5$\,mm), proximity to primary drainage canals ($<3$\,km), and extended dry spells ($>14$\,days) strongly increase the log-odds of severe Fire Radiative Power ($\text{FRP} \ge 12.0$\,MW).}
\label{fig:shap_importance}
\end{figure}

The TreeSHAP beeswarm distribution revealed three critical non-linear interaction thresholds governing tropical peat fire severity:
\begin{enumerate}
    \item \textbf{Antecedent Rainfall Deficit Threshold}: 7-day cumulative precipitation ($P_{\text{7d}}$) emerged as the single most dominant predictor ($\text{mean } |\text{SHAP}| = 0.0329$). High rainfall values (red points) exerted strong negative attributions, suppressing fire intensity. However, when cumulative rainfall dropped below 5.0\,mm, SHAP values exhibited a sharp positive inflection, increasing severe fire log-odds by up to $+0.18$.
    \item \textbf{Canal Drainage Proximity Corridor}: In the ex-PLG landscape, proximity to primary drainage canals ($d_{\text{canal}}$) exerted the second-highest attribution (0.0322). Within a 3,000-meter corridor from primary canal axes, SHAP values rose precipitously, confirming that artificial drainage channels create dry micro-climatic zones where fire energy release is maximized.
    \item \textbf{The 14-Day Dry Spell Tipping Point}: Dry spell duration ($D_{\text{spell}}$) showed negligible positive attribution during the initial 7 days of dry weather. However, between 12 and 14 consecutive rainless days, SHAP values surged exponentially as phreatic water tables dropped below the upper fibrous peat zone, transitioning the ecosystem into extreme combustion risk.
\end{enumerate}
